\documentclass[aspectratio=169,10pt]{beamer}
\usepackage[T1]{fontenc}
\usepackage[utf8]{inputenc}
\usepackage[spanish,es-tabla]{babel}
\usepackage{lmodern}
\usepackage{helvet}
\usepackage{amsmath,booktabs}
\usetheme{P2}
% Criterio de servicio para el reentreno restringido.
% Sin resultados pendientes: solo regla de aceptacion y limite de interpretacion.
\newcommand{\rarEstado}{Criterio de servicio}
\newcommand{\rarDetalle}{Ventanas disjuntas por d\'ias; la ventana anual ya estaba inspeccionada.}
\newcommand{\rarQuiebres}{menos quiebres regionales}
\newcommand{\rarUnidades}{m\'as ventas, ninguna tienda peor}
\newcommand{\rarRetorno}{m\'as retorno que la ingenua}

\title{\textquestiondown C\'omo repartir el inventario entre tres tiendas?}
\subtitle{Precios e inventario con aprendizaje por refuerzo}
\author{Leonardo Candio \and Dimael Rivas}
\date{DS5345 \textperiodcentered\ Fase 2 \textperiodcentered\ V\'ia B}

\begin{document}

% ---------- 1. Portada ----------
\begin{frame}[plain,noframenumbering]
\vspace{0.5cm}
\pill{DS5345 \textperiodcentered\ Fase 2 \textperiodcentered\ V\'ia B}
\hspace{0.2cm}\pill[p2ambersoft]{365 d\'ias \textperiodcentered\ 3 tiendas}

\vspace{0.3cm}
{\fontsize{34}{36}\selectfont\bfseries \textquestiondown C\'omo repartir\\[-2pt] el inventario entre tres tiendas?\par}
\vspace{0.15cm}
{\large\color{p2muted} Precios e inventario con aprendizaje por refuerzo\par}

\vspace{0.4cm}
\pill{Tienda 1}\hspace{0.15cm}\pill{Tienda 2}\hspace{0.15cm}\pill{Tienda 3}
\hspace{0.5cm}{\small\color{p2muted}reponen desde un almac\'en regional compartido}

\vspace{0.45cm}
{\small Leonardo Candio \quad\textcolor{p2line}{$|$}\quad Dimael Rivas}

\vspace{0.15cm}
{\scriptsize\color{p2teal} Simulaci\'on con datos de ventas y precios M5\par}
\end{frame}

\section{Tres tiendas, un almac\'en}

% ---------- 2. Portadilla bloque 1 ----------
\begin{frame}[plain]
\vspace{1.1cm}
{\fontsize{52}{52}\selectfont\bfseries\color{p2teal} I\par}
\vspace{0.1cm}
{\fontsize{24}{26}\selectfont\bfseries Tres tiendas, un almac\'en\par}
\vspace{0.15cm}
\pill[p2ambersoft]{Bloque 1 \textperiodcentered\ el problema}
\hspace{0.3cm}{\small\color{p2muted}pedido diario \textperiodcentered\ tres formas de pedir \textperiodcentered\ m\'etodos}
\end{frame}

% ---------- 3. Tres formas de pedir ----------
\begin{frame}
\frametitle{Tres formas de pedir}
\framesubtitle{Men\'u discreto: multiplicador de precio $\times$ cobertura de pedido}
\begin{columns}[T,onlytextwidth]
\column{0.31\textwidth}
\begin{block}{Ingenua}
\scriptsize Repite las ventas de hace una semana, aproximadas al men\'u. No aprende.
\end{block}
\column{0.31\textwidth}
\begin{block}{Ego\'ista}
\scriptsize Aprende a mejorar el resultado de su propia tienda.
\end{block}
\column{0.31\textwidth}
\begin{block}{Cooperativa}
\scriptsize Aprende usando el resultado de las tres y el estado del recurso compartido.
\end{block}
\end{columns}
\vspace{0.3cm}
\begin{center}
{\scriptsize\color{p2muted} La ego\'ista y la cooperativa aprenden por refuerzo (Q-learning tabular y PPO): prueban pedidos y se quedan con lo que da mejor resultado.\par}
\end{center}
\end{frame}

% ---------- 4. Secuencia diaria ----------
\begin{frame}
\frametitle{Cada d\'ia, las tiendas consultan el stock y hacen su pedido}
\framesubtitle{Si piden de m\'as, el almac\'en raciona en proporci\'on}
\begin{center}
\begin{tikzpicture}[
  box/.style={rounded corners=6pt,fill=white,draw=p2line,inner sep=8pt,align=center,font=\small},
  arr/.style={-{Stealth[length=7pt]},thick,color=p2teal}]
  \node[box] (v) {ventas\\pasadas};
  \node[box,fill=p2tealsoft,right=1.2cm of v,minimum width=2.4cm] (p) {\textbf{pedido}\\precio $\times$ cobertura};
  \node[box,fill=p2ambersoft,right=1.2cm of p,minimum width=2.4cm] (r) {\textbf{reparto}\\proporcional};
  \node[box,right=1.2cm of r,minimum width=2.4cm] (l) {llega\\al d\'ia siguiente};
  \draw[arr] (v) -- (p);
  \draw[arr] (p) -- (r);
  \draw[arr] (r) -- (l);
\end{tikzpicture}
\end{center}
\vspace{0.2cm}
\begin{center}
{\scriptsize\color{p2muted} Esquema de ejemplo con pedidos ilustrativos.\par}
\end{center}
\end{frame}

% ---------- 5. M5 vs simulado ----------
\begin{frame}
\frametitle{M5 real vs.\ simulado}
\framesubtitle{Las tres pol\'iticas usan el mismo precio de cada d\'ia}
\begin{columns}[T,onlytextwidth]
\column{0.47\textwidth}
\begin{center}\pill{M5 aporta}\end{center}
\vspace{0.1cm}
\begin{itemize}
  \item Precios y calendario
  \item Ventas $\approx$ demanda de referencia
\end{itemize}
\column{0.47\textwidth}
\begin{center}\pill[p2ambersoft]{El simulador supone}\end{center}
\vspace{0.1cm}
\begin{itemize}
  \item Pedidos, stock, faltantes
  \item Costes y elasticidad con redondeo
\end{itemize}
\end{columns}
\vspace{0.35cm}
\begin{center}
{\scriptsize\color{p2teal} Simulaci\'on con datos de ventas y precios M5\par}
\end{center}
\end{frame}

% ---------- 6. Escalera ----------
\begin{frame}
\frametitle{M\'etodos de aprendizaje}
\framesubtitle{Dos herramientas de entrenamiento}
\begin{columns}[T,onlytextwidth]
\column{0.40\textwidth}
\begin{center}
\begin{tikzpicture}[
  paso/.style={rounded corners=6pt,fill=white,draw=p2teal,thick,inner sep=7pt,minimum width=3.6cm,align=center,font=\small\bfseries},
  arr/.style={-{Stealth[length=6pt]},thick,color=p2teal}]
  \node[paso] (t) {Tabular};
  \node[paso,below=0.35cm of t] (d) {Double DQN};
  \node[paso,below=0.35cm of d] (p) {PPO};
  \draw[arr] (t) -- (d);
  \draw[arr] (d) -- (p);
\end{tikzpicture}
\end{center}
\column{0.56\textwidth}
\begin{block}{PBRS \normalfont\scriptsize\color{p2muted} (ajuste de recompensa con un potencial)}
\scriptsize $F_t = \gamma\Phi(s_{t+1}) - \Phi(s_t)$: conserva el orden de acciones. El potencial terminal es cero.
\end{block}
\vspace{0.15cm}
\begin{block}{Dyna \normalfont\scriptsize\color{p2muted} (modelo que genera transiciones)}
\scriptsize Usa un modelo del entorno para generar experiencias de entrenamiento.
\end{block}
\vspace{0.15cm}
{\scriptsize\color{p2muted} Q-learning y Double DQN estiman el valor de las acciones. PPO ajusta la pol\'itica.\par}
\end{columns}
\end{frame}

\section{Qu\'e cambi\'o en el simulador}

% ---------- 7. Portadilla bloque 2 ----------
\begin{frame}[plain]
\vspace{0.9cm}
{\fontsize{52}{52}\selectfont\bfseries\color{p2amber} II\par}
\vspace{0.1cm}
{\fontsize{24}{26}\selectfont\bfseries Qu\'e cambi\'o en el simulador\par}
\vspace{0.15cm}
\pill{Las tres pol\'iticas enfrentan la misma demanda}
\vspace{0.25cm}
{\small\color{p2muted} un producto \textperiodcentered\ 365 d\'ias \textperiodcentered\ 3 tiendas \textperiodcentered\ mismo precio del d\'ia \textperiodcentered\ tres entrenamientos\par}
\vspace{0.15cm}
{\small\color{p2ink} Un quiebre es un d\'ia en que una tienda no puede atender toda su demanda.\par}
\end{frame}

% ---------- 8. Incidencia ----------
\begin{frame}
\frametitle{El m\'etodo cooperativo tuvo menos d\'ias con quiebre}
\framesubtitle{En las tres corridas: 88--91\% menos frente a la ingenua; 34--69\% menos frente a la ego\'ista}
\begin{columns}[T,onlytextwidth]
\column{0.46\textwidth}
\begin{center}
\begin{tikzpicture}
\begin{axis}[
  width=5.8cm,height=3.4cm,
  ybar, bar width=9pt,
  symbolic x coords={C1,C2,C3,E1,E2,E3,N},
  xtick=data,
  xticklabels={Coop-1,Coop-2,Coop-3,Ego-1,Ego-2,Ego-3,Ing},
  x tick label style={font=\tiny},
  ylabel={d\'ias-tienda con quiebre}, ylabel style={font=\scriptsize},
  nodes near coords, nodes near coords style={font=\tiny\bfseries},
  ymin=0, ymax=1120, enlarge x limits=0.12,
  axis lines*=left, tick style={draw=none}, axis line style={color=p2line},
]
\addplot[fill=p2teal,draw=p2teal] coordinates {(C1,116) (C2,122) (C3,93)};
\addplot[fill=p2amber,draw=p2amber] coordinates {(E1,376) (E2,185) (E3,183)};
\addplot[fill=p2muted,draw=p2muted] coordinates {(N,1015)};
\end{axis}
\end{tikzpicture}
\end{center}
\column{0.50\textwidth}
\begin{center}
{\Large\bfseries 88--91\%} menos que la ingenua\quad {\Large\bfseries 34--69\%} menos que la ego\'ista
\end{center}
\vspace{0.15cm}
\begin{itemize}
  \item Tasa media de d\'ias con quiebre: ingenua 93\%, ego\'ista 23\%, cooperativa 10\%
  \item 1\,092 d\'ias-tienda elegibles por brazo
\end{itemize}
\end{columns}
\vspace{0.08cm}
{\scriptsize\color{p2muted} Ingenua $=$ heur\'istica de ventas rezagadas; referencia principal. Ego\'ista $=$ emparejada por corrida.\par}
\end{frame}

% ---------- 9. Ventas ----------
\begin{frame}
\frametitle{Menos quiebres no siempre signific\'o m\'as ventas}
\framesubtitle{Ventas con la misma demanda (57\,184 u) para las tres pol\'iticas}
\begin{columns}[T,onlytextwidth]
\column{0.46\textwidth}
\begin{center}
\begin{tikzpicture}
\begin{axis}[
  width=5.8cm,height=3.4cm,
  ybar, bar width=11pt,
  symbolic x coords={R1,R2,R3},
  xtick=data,
  xticklabels={Corrida 1,Corrida 2,Corrida 3},
  x tick label style={font=\tiny},
  ylabel={unidades vendidas}, ylabel style={font=\scriptsize},
  nodes near coords, nodes near coords style={font=\tiny\bfseries},
  ymin=40000, ymax=54500, enlarge x limits=0.25,
  axis lines*=left, tick style={draw=none}, axis line style={color=p2line},
]
\addplot[fill=p2amber,draw=p2amber] coordinates {(R1,41747) (R2,52590) (R3,52360)};
\addplot[fill=p2teal,draw=p2teal] coordinates {(R1,52133) (R2,52473) (R3,52008)};
\end{axis}
\end{tikzpicture}
\end{center}
\column{0.50\textwidth}
\begin{itemize}
  \item Corrida 3: ego\'ista 52\,360 vs.\ cooperativa 52\,008
  \item Solo 1 de 3 supera a la ego\'ista en ventas y retorno
  \item 3 de 3 supera a la ingenua en ambas
\end{itemize}
\vspace{0.1cm}
\begin{center}
\pill{Cooperativa}\hspace{0.1cm}\pill[p2ambersoft]{Ego\'ista}
\end{center}
\end{columns}
\vspace{0.06cm}
{\scriptsize\color{p2muted} Con la misma demanda, cada unidad faltante es una venta perdida.\par}
\end{frame}

% ---------- 10. Demo ----------
\begin{frame}
\frametitle{Un a\~no en 36 segundos}
\framesubtitle{\texttt{video.mp4} en el mismo ZIP \textperiodcentered\ comparar Ingenua, Ego\'ista y Cooperativa}
\begin{columns}[T,onlytextwidth]
\column{0.52\textwidth}
\begin{center}
\begin{tikzpicture}[
  cell/.style={rounded corners=3pt,minimum width=0.85cm,minimum height=0.62cm,inner sep=1pt,font=\tiny},
  rowl/.style={font=\scriptsize\bfseries\color{p2ink},anchor=east}]
  \foreach \r/\name/\col in {0/Cooperativa/p2tealsoft,1/Ego\'ista/p2ambersoft,2/Ingenua/p2paper} {
    \node[rowl] at (-0.2,-\r*0.85) {\name};
    \foreach \cc in {0,1,2} {
      \node[cell,fill=\col,draw=p2line] at (\cc*1.0,-\r*0.85) {T\_\the\numexpr\cc+1\relax};
    }
  }
  \node[font=\scriptsize\color{p2muted}] at (1.0,0.65) {d\'ia 1 $\rightarrow$ 365 sincronizado};
\end{tikzpicture}
\end{center}
\column{0.44\textwidth}
\begin{itemize}
  \item 3 filas $\times$ 3 tiendas, mismo d\'ia
  \item Color $=$ unidades perdidas, no solo si hubo quiebre
  \item Quiebres, faltantes y retorno acumulados
  \item Una corrida de la simulaci\'on
\end{itemize}
\end{columns}
\end{frame}

% ---------- 11. Cierre ----------
\begin{frame}
\frametitle{Para evaluar la reposici\'on, contamos los d\'ias de quiebre y las unidades perdidas}
\begin{center}
\vspace{0.5cm}
\begin{tikzpicture}
\node[rounded corners=8pt,fill=white,draw=p2teal,thick,inner sep=14pt,align=center] {%
{\large\bfseries\color{p2ink} D\'ias con quiebre \textcolor{p2line}{$+$} unidades perdidas}\\[4pt]
{\small\color{p2muted} las dos m\'etricas, en las tres corridas}};
\end{tikzpicture}
\vspace{0.5cm}

{\small\color{p2muted} Gracias \textperiodcentered\ preguntas\par}
\end{center}
\end{frame}

\end{document}
